\chapter{Clase 21}

\section{Hidrodinámica}

La hidrodinámica describe el movimiento de los fluidos. Para ello, primero se especifica la velocidad en cada punto y se distinguen las curvas que representan instantáneamente el campo de las trayectorias de sus partículas. Finalmente, se formula la conservación de la masa.

\subsection{Campo de velocidad y momento lineal}

El movimiento de un fluido continuo se describe mediante el campo de velocidad
\[
\vec v=\vec v(\vec r,t),
\]
que asigna a cada posición $\vec r$ y tiempo $t$ la velocidad de la porción de fluido que ocupa ese punto. En coordenadas cartesianas,
\[
\vec v(\vec r,t)=v_x(\vec r,t)\,\hat e_x+v_y(\vec r,t)\,\hat e_y+v_z(\vec r,t)\,\hat e_z.
\]
La velocidad es un campo euleriano: se describe cómo cambia en posiciones fijas del espacio, aunque las partículas de fluido se muevan.

Una porción infinitesimal de masa $dm=\rho\,dV$, con densidad $\rho$, tiene momento lineal
\[
d\vec p=\vec v\,dm=\rho\vec v\,dV.
\]
Por tanto, la densidad de momento lineal (momento por unidad de volumen) es $\rho\vec v$.

\subsection{Visualización del flujo}

Un campo de velocidad puede representarse mediante flechas que muestran la dirección y magnitud de $\vec v$ en distintos puntos. También se utilizan curvas geométricas que permiten visualizar cómo se orienta el flujo.

\textbf{Línea de corriente.} En un instante fijo $t_0$, una línea de corriente es una curva cuya tangente en cada punto es paralela a la velocidad local. Si $\vec r_s(s)$ la parametriza, entonces
\[
\frac{d\vec r_s}{ds}=\alpha(s)\,\vec v\bigl(\vec r_s(s),t_0\bigr),
\]
donde $\alpha(s)$ es un factor escalar que depende de la parametrización. En componentes, cuando las velocidades correspondientes no son nulas, esto equivale a
\[
\frac{dx}{v_x}=\frac{dy}{v_y}=\frac{dz}{v_z}.
\]
La ecuación se resuelve manteniendo fijo el tiempo $t_0$: cada línea describe la dirección instantánea del flujo, no necesariamente el recorrido posterior de una partícula.

En un punto donde $\vec v\ne0$ solo hay una dirección de velocidad; por eso, dos líneas de corriente no se cruzan. En los puntos de estancamiento, donde $\vec v=0$, la dirección tangente no queda definida.

Algunos ejemplos sencillos son:
\begin{itemize}
  \item Si $\vec v=U_0\hat e_x$, con $U_0$ constante, las líneas de corriente son rectas paralelas al eje $x$.
  \item Si $\vec v=\vec\Omega\times\vec r$ describe una rotación rígida alrededor de un eje, las líneas de corriente son circunferencias centradas en dicho eje.
\end{itemize}

\textbf{Tubo de corriente.} Es la región limitada por un conjunto de líneas de corriente. Como la velocidad es tangente a sus paredes, el fluido no las atraviesa: la masa solo entra o sale por las secciones de sus extremos.

\begin{figure}[H]
\centering
\begin{tikzpicture}[
  >=Latex,
  line cap=round,
  line join=round,
  font=\small,
  flujo/.style={->, blue!75!black, line width=1.15pt},
  linea/.style={blue!55!black, line width=0.9pt}
]
\node[font=\small\bfseries] at (-4.45,1.8) {campo uniforme};
\foreach \y in {-1.05,-0.35,0.35,1.05} {
  \draw[linea] (-5.9,\y) -- (-3.0,\y);
  \draw[flujo] (-4.5,\y) -- (-3.7,\y);
}

\node[font=\small\bfseries] at (-0.2,1.8) {rotación rígida};
\fill[black!70] (-0.2,0) circle (1.5pt);
\foreach \radio in {0.55,1.0,1.45} {
  \draw[linea] (-0.2,0) circle (\radio);
  \draw[flujo] ({-0.2+\radio},0)
    arc[start angle=0,end angle=42,radius=\radio];
}
\node[fill=white, inner sep=1pt] at (-0.2,-1.8) {$\vec v=\vec\Omega\times\vec r$};

\node[font=\small\bfseries] at (4.35,1.8) {tubo de corriente};
\fill[blue!5]
  (2.15,1.1) .. controls (3.5,1.05) and (4.8,0.72) .. (6.25,0.62)
  -- (6.25,-0.62) .. controls (4.8,-0.72) and (3.5,-1.05) .. (2.15,-1.1)
  -- cycle;
\draw[linea] (2.15,1.1) .. controls (3.5,1.05) and (4.8,0.72) .. (6.25,0.62);
\draw[linea] (2.15,-1.1) .. controls (3.5,-1.05) and (4.8,-0.72) .. (6.25,-0.62);
\draw[densely dashed, gray!75] (2.3,-1.06) -- (2.3,1.06);
\draw[densely dashed, gray!75] (6.1,-0.63) -- (6.1,0.63);
\node[above, fill=white, inner sep=1pt] at (2.3,1.1) {$A_1$};
\node[above, fill=white, inner sep=1pt] at (6.1,0.63) {$A_2$};
\draw[flujo] (2.8,0.48) -- (3.8,0.45);
\draw[flujo] (3.8,-0.05) -- (4.8,-0.05);
\draw[flujo] (4.7,-0.43) -- (5.6,-0.40);
\end{tikzpicture}
\par\medskip
\textbf{Figura 34.} Representaciones de un flujo uniforme, de una rotación rígida y de un tubo de corriente. Las paredes del tubo siguen líneas de corriente, por lo que no hay flujo de masa a través de ellas.
\end{figure}

\subsection{Trayectorias de las partículas y líneas de corriente}

La \textbf{trayectoria} (o \emph{pathline}) es el camino que recorre una partícula concreta a medida que transcurre el tiempo. Si su posición es $\vec r_p(t)$, satisface la ecuación de movimiento
\[
\frac{d\vec r_p}{dt}=\vec v\bigl(\vec r_p(t),t\bigr),
\qquad
\vec r_p(t_0)=\vec r_0.
\]
A diferencia de la línea de corriente, que se construye para un instante fijo, la trayectoria sigue a la misma partícula y, por tanto, incorpora la evolución temporal del campo de velocidad.

En un flujo estacionario, el campo de velocidad no depende explícitamente del tiempo y las trayectorias coinciden con las líneas de corriente. En un flujo no estacionario, en general son distintas. Por ejemplo, para el campo plano $\vec v=(U,\gamma t)$, con $U\ne0$ y $\gamma$ constantes, las líneas de corriente en un instante dado son rectas de pendiente $\gamma t/U$. En cambio, una partícula que parte de $(x_0,y_0)$ en el instante $t_0$ sigue
\[
x(t)=x_0+U(t-t_0),
\qquad
y(t)=y_0+\frac{\gamma}{2}(t^2-t_0^2),
\]
que es una trayectoria parabólica si $\gamma\ne0$.

\subsection{Conservación de la masa}

La masa de un sistema material se conserva. Si $V_m(t)$ es un volumen que se mueve junto con las mismas partículas de fluido, su masa es
\[
M(t)=\int_{V_m(t)}\rho\,dV,
\qquad
\frac{d}{dt}\int_{V_m(t)}\rho\,dV=0.
\]

También se puede describir el balance en un volumen de control fijo $V$, limitado por una superficie cerrada $S$ con normal exterior $\hat n$. La variación de la masa dentro de $V$ más el flujo neto de masa que sale a través de $S$ debe ser cero:
\[
\frac{d}{dt}\int_V\rho\,dV
+\oint_S\rho\,\vec v\cdot\hat n\,dS=0.
\]
El primer término mide la acumulación de masa dentro del volumen; el segundo es el caudal másico neto saliente. Aplicando el teorema de Gauss y considerando volúmenes arbitrarios, se obtiene la forma local de la ecuación de continuidad:
\[
\boxed{\frac{\partial\rho}{\partial t}+\nabla\cdot(\rho\vec v)=0.}
\]
En términos de la derivada material,
\[
\frac{D\rho}{Dt}
=\frac{\partial\rho}{\partial t}+\vec v\cdot\nabla\rho,
\qquad
\frac{D\rho}{Dt}+\rho\,\nabla\cdot\vec v=0.
\]
Así, un fluido de densidad constante satisface $\nabla\cdot\vec v=0$; se dice entonces que su movimiento es incompresible.

En régimen estacionario, la masa que atraviesa por unidad de tiempo cualquier sección transversal de un tubo de corriente es la misma:
\[
\dot m=\int_A\rho\,\vec v\cdot\hat n\,dA=\text{constante}.
\]
Si la densidad y la velocidad normal son uniformes en cada sección, esto se reduce a
\[
\rho_1A_1v_1=\rho_2A_2v_2.
\]
Para un fluido incompresible, $\rho$ es constante y por tanto $A_1v_1=A_2v_2$: al disminuir el área de la sección, aumenta la rapidez media del flujo.